\documentclass[12pt,a4paper]{article}

% Packages
\usepackage{geometry}
\usepackage{graphicx}
\usepackage{booktabs}
\usepackage{longtable}
\usepackage{array}
\usepackage{multirow}
\usepackage{colortbl}
\usepackage{xcolor}
\usepackage{float}
\usepackage{amsmath}
\usepackage{amssymb}
\usepackage{caption}
\usepackage{subcaption}
\usepackage{hyperref}

% Page settings
\geometry{left=2.5cm,right=2.5cm,top=2.5cm,bottom=2.5cm}

% Hyperlink settings
\hypersetup{
    colorlinks=true,
    linkcolor=blue,
    filecolor=magenta,
    urlcolor=cyan,
    citecolor=green,
}

% Custom colors
\definecolor{lightgray}{gray}{0.9}
\definecolor{excellent}{RGB}{46, 139, 87}
\definecolor{good}{RGB}{60, 179, 113}
\definecolor{fair}{RGB}{255, 165, 0}
\definecolor{poor}{RGB}{220, 20, 60}

% Table styles
\newcolumntype{L}[1]{>{\raggedright\arraybackslash}m{#1}}
\newcolumntype{C}[1]{>{\centering\arraybackslash}m{#1}}
\newcolumntype{R}[1]{>{\raggedleft\arraybackslash}m{#1}}

% Title information
\title{\textbf{Statistical Reliability Analysis of LLM Multi-Run Evaluations}\\
\large Assessing the Consistency and Stability of Repeated AI Model Runs}
\author{LLM Consensus Analysis Framework}
\date{\today}

\begin{document}

\maketitle

\tableofcontents
\newpage


\section{Background and Objectives}

\subsection{Research Background}
In materials science decision-making, Large Language Models (LLMs) are widely used for evaluating and screening candidate materials. However, whether the same LLM model can produce consistent and reliable results across multiple runs is a critical statistical question. This study aims to systematically evaluate the reliability of LLM models in multiple independent runs through statistical analysis.

\subsection{Research Objectives}
The core objective of this study is to test the following hypothesis:
\begin{quote}
\textit{Multiple independent runs of the same AI model are statistically reliable, with consistent scoring and stable decision-making.}
\end{quote}

Specifically, we evaluate reliability across the following dimensions:
\begin{itemize}
    \item \textbf{Scoring Consistency}: Variation in scores for the same material across different runs
    \item \textbf{Inter-rater Reliability}: Correlation of scores between different runs
    \item \textbf{Decision Stability}: Consistency of final winner selection
\end{itemize}

\subsection{Data Source}
This study analyzes four mainstream LLM models (GPT-5, Grok-4, Claude Opus 4.5, Gemini 3 Pro) across 11 independent runs. Each run evaluates 10 candidate materials across 6 dimensions (0-10 points) and selects the optimal material.

\newpage


\section{Methodology}

\subsection{Statistical Metrics}

\subsubsection{Coefficient of Variation (CV)}
Coefficient of variation measures the relative dispersion of scores:
\begin{equation}
    CV = \frac{\sigma}{\mu} \times 100\%
\end{equation}
where $\sigma$ is standard deviation and $\mu$ is the mean. Lower CV values indicate higher scoring consistency.

\subsubsection{Intraclass Correlation Coefficient (ICC)}
ICC measures the consistency between different runs (raters). Based on the ANOVA method by Shrout \& Fleiss (1979), we calculate ICC(3,1) for a two-way fixed-effects model with absolute agreement:
\begin{equation}
    ICC(3,1) = \frac{MS_R - MS_E}{MS_R + (k-1)MS_E}
\end{equation}
where $MS_R$ is mean square between targets (materials), $MS_E$ is mean square error, and $k$ is the number of raters (runs). ICC values closer to 1 indicate higher consistency.

\subsubsection{Consistency Ratio}
Consistency ratio is defined as the ratio of score range to mean:
\begin{equation}
    CR = \frac{Range}{Mean}
\end{equation}
This metric reflects the relative fluctuation range of scores.

\subsubsection{Decision Consistency Rate}
Decision consistency rate is the proportion of runs selecting the same winner:
\begin{equation}
    DC = \frac{N_{most\_frequent}}{N_{total}} \times 100\%
\end{equation}

\subsubsection{Information Entropy}
Information entropy measures the uncertainty in winner distribution:
\begin{equation}
    H = -\sum_{i} p_i \log_2(p_i)
\end{equation}
where $p_i$ is the probability of formula $i$ being selected as the winner. Lower entropy values indicate more concentrated decision-making.

\subsection{Reliability Assessment Criteria}
We adopt the following reliability assessment standards based on statistical metric ranges:

\begin{table}[H]
\centering
\caption{Reliability Assessment Criteria}
\begin{tabular}{lll}
\toprule
\textbf{Metric} & \textbf{Grade} & \textbf{Threshold Range} \\
\midrule
\multirow{3}{*}{Coefficient of Variation (CV)} & Excellent & CV $<$ 10\% \\
& Good & 10\% $\le$ CV $<$ 20\% \\
& Fair & CV $\ge$ 20\% \\
\midrule
\multirow{3}{*}{Intraclass Correlation (ICC)} & Excellent & ICC $>$ 0.8 \\
& Good & 0.6 $<$ ICC $\le$ 0.8 \\
& Fair & ICC $\le$ 0.6 \\
\midrule
\multirow{3}{*}{Decision Consistency Rate (DC)} & Excellent & DC $\ge$ 80\% \\
& Good & 60\% $\le$ DC $<$ 80\% \\
& Fair & DC $<$ 60\% \\
\bottomrule
\end{tabular}
\end{table}

\newpage


\section{gpt-5 Reliability Analysis}

\subsection{Scoring Consistency Analysis}

The following table displays the coefficient of variation (CV) across scoring dimensions for gpt-5:

\begin{table}[H]
\centering
\caption{gpt-5 - Scoring Variation Coefficient Analysis}
\begin{tabular}{llccc}
\toprule
\textbf{Scoring Dimension} & \textbf{Mean CV(\%)} & \textbf{Max CV(\%)} & \textbf{Min CV(\%)} & \textbf{Consistency Ratio} \\
\midrule
Mechanical\_Safety & 6.50 & 9.93 & 4.59 & 0.157 \\
Swelling\_Performance & 9.58 & 14.46 & 5.84 & 0.284 \\
Endothelialization & 8.07 & 14.35 & 3.73 & 0.246 \\
SMC\_inhibition & 10.46 & 16.86 & 5.32 & 0.310 \\
Anti\_inflammation & 10.99 & 20.73 & 5.58 & 0.357 \\
Thrombogenicity & 11.33 & 19.13 & 4.59 & 0.323 \\
Total\_Score & 5.80 & 8.92 & 3.65 & 0.189 \\
\bottomrule
\end{tabular}
\end{table}

\subsection{Inter-rater Reliability Analysis}

The intraclass correlation coefficient (ICC) reflects scoring consistency across different runs:

\begin{table}[H]
\centering
\caption{gpt-5 - Intraclass Correlation Coefficient}
\begin{tabular}{lc}
\toprule
\textbf{Scoring Dimension} & \textbf{ICC} \\
\midrule
Mechanical\_Safety & 0.7072 \\
Swelling\_Performance & 0.6563 \\
Endothelialization & 0.8912 \\
SMC\_inhibition & 0.8356 \\
Anti\_inflammation & 0.8338 \\
Thrombogenicity & 0.8864 \\
Total\_Score & 0.8776 \\
\bottomrule
\end{tabular}
\end{table}

\subsection{Decision Consistency Analysis}


\begin{itemize}
    \item \textbf{Total Runs}: 11 runs
    \item \textbf{Most Common Winner}: Formula 5
    \item \textbf{Occurrence Count}: 10 times
    \item \textbf{Decision Consistency Rate}: 90.9\%
    \item \textbf{Information Entropy}: 0.439 (Max Entropy: 1.000)
\end{itemize}

Winner formula distribution:

\begin{table}[H]
\centering
\caption{gpt-5 - Winner Formula Distribution}
\begin{tabular}{cc}
\toprule
\textbf{Formula} & \textbf{Count} \\
\midrule
Formula 10 & 1 \\
Formula 5 & 10 \\
\bottomrule
\end{tabular}
\end{table}


\newpage


\section{grok-4 Reliability Analysis}

\subsection{Scoring Consistency Analysis}

The following table displays the coefficient of variation (CV) across scoring dimensions for grok-4:

\begin{table}[H]
\centering
\caption{grok-4 - Scoring Variation Coefficient Analysis}
\begin{tabular}{llccc}
\toprule
\textbf{Scoring Dimension} & \textbf{Mean CV(\%)} & \textbf{Max CV(\%)} & \textbf{Min CV(\%)} & \textbf{Consistency Ratio} \\
\midrule
Mechanical\_Safety & 6.32 & 10.17 & 4.97 & 0.166 \\
Swelling\_Performance & 9.22 & 15.81 & 3.38 & 0.284 \\
Endothelialization & 9.58 & 14.27 & 0.00 & 0.288 \\
SMC\_inhibition & 14.47 & 23.36 & 10.31 & 0.416 \\
Anti\_inflammation & 10.38 & 21.08 & 0.00 & 0.297 \\
Thrombogenicity & 12.23 & 37.90 & 4.41 & 0.369 \\
Total\_Score & 5.86 & 11.37 & 3.17 & 0.188 \\
\bottomrule
\end{tabular}
\end{table}

\subsection{Inter-rater Reliability Analysis}

The intraclass correlation coefficient (ICC) reflects scoring consistency across different runs:

\begin{table}[H]
\centering
\caption{grok-4 - Intraclass Correlation Coefficient}
\begin{tabular}{lc}
\toprule
\textbf{Scoring Dimension} & \textbf{ICC} \\
\midrule
Mechanical\_Safety & 0.4084 \\
Swelling\_Performance & 0.2277 \\
Endothelialization & 0.8598 \\
SMC\_inhibition & 0.2300 \\
Anti\_inflammation & 0.7017 \\
Thrombogenicity & 0.6999 \\
Total\_Score & 0.6300 \\
\bottomrule
\end{tabular}
\end{table}

\subsection{Decision Consistency Analysis}


\begin{itemize}
    \item \textbf{Total Runs}: 11 runs
    \item \textbf{Most Common Winner}: Formula 4
    \item \textbf{Occurrence Count}: 7 times
    \item \textbf{Decision Consistency Rate}: 63.6\%
    \item \textbf{Information Entropy}: 1.241 (Max Entropy: 1.585)
\end{itemize}

Winner formula distribution:

\begin{table}[H]
\centering
\caption{grok-4 - Winner Formula Distribution}
\begin{tabular}{cc}
\toprule
\textbf{Formula} & \textbf{Count} \\
\midrule
Formula 3 & 1 \\
Formula 4 & 7 \\
Formula 5 & 3 \\
\bottomrule
\end{tabular}
\end{table}


\newpage


\section{claude-opus-4-5-20251101 Reliability Analysis}

\subsection{Scoring Consistency Analysis}

The following table displays the coefficient of variation (CV) across scoring dimensions for claude-opus-4-5-20251101:

\begin{table}[H]
\centering
\caption{claude-opus-4-5-20251101 - Scoring Variation Coefficient Analysis}
\begin{tabular}{llccc}
\toprule
\textbf{Scoring Dimension} & \textbf{Mean CV(\%)} & \textbf{Max CV(\%)} & \textbf{Min CV(\%)} & \textbf{Consistency Ratio} \\
\midrule
Mechanical\_Safety & 3.54 & 6.54 & 0.00 & 0.084 \\
Swelling\_Performance & 7.73 & 20.55 & 0.00 & 0.194 \\
Endothelialization & 4.65 & 10.36 & 0.00 & 0.113 \\
SMC\_inhibition & 7.89 & 12.53 & 0.00 & 0.221 \\
Anti\_inflammation & 8.75 & 14.14 & 4.25 & 0.185 \\
Thrombogenicity & 8.08 & 17.13 & 0.00 & 0.218 \\
Total\_Score & 4.27 & 6.12 & 3.06 & 0.127 \\
\bottomrule
\end{tabular}
\end{table}

\subsection{Inter-rater Reliability Analysis}

The intraclass correlation coefficient (ICC) reflects scoring consistency across different runs:

\begin{table}[H]
\centering
\caption{claude-opus-4-5-20251101 - Intraclass Correlation Coefficient}
\begin{tabular}{lc}
\toprule
\textbf{Scoring Dimension} & \textbf{ICC} \\
\midrule
Mechanical\_Safety & 0.8970 \\
Swelling\_Performance & 0.7807 \\
Endothelialization & 0.9663 \\
SMC\_inhibition & 0.6897 \\
Anti\_inflammation & 0.8188 \\
Thrombogenicity & 0.7797 \\
Total\_Score & 0.8708 \\
\bottomrule
\end{tabular}
\end{table}

\subsection{Decision Consistency Analysis}


\begin{itemize}
    \item \textbf{Total Runs}: 11 runs
    \item \textbf{Most Common Winner}: Formula 4
    \item \textbf{Occurrence Count}: 7 times
    \item \textbf{Decision Consistency Rate}: 63.6\%
    \item \textbf{Information Entropy}: 1.309 (Max Entropy: 1.585)
\end{itemize}

Winner formula distribution:

\begin{table}[H]
\centering
\caption{claude-opus-4-5-20251101 - Winner Formula Distribution}
\begin{tabular}{cc}
\toprule
\textbf{Formula} & \textbf{Count} \\
\midrule
Formula 3 & 2 \\
Formula 4 & 7 \\
Formula 5 & 2 \\
\bottomrule
\end{tabular}
\end{table}


\newpage


\section{gemini-3-pro-preview Reliability Analysis}

\subsection{Scoring Consistency Analysis}

The following table displays the coefficient of variation (CV) across scoring dimensions for gemini-3-pro-preview:

\begin{table}[H]
\centering
\caption{gemini-3-pro-preview - Scoring Variation Coefficient Analysis}
\begin{tabular}{llccc}
\toprule
\textbf{Scoring Dimension} & \textbf{Mean CV(\%)} & \textbf{Max CV(\%)} & \textbf{Min CV(\%)} & \textbf{Consistency Ratio} \\
\midrule
Mechanical\_Safety & 10.28 & 20.55 & 0.00 & 0.220 \\
Swelling\_Performance & 7.81 & 26.35 & 0.00 & 0.167 \\
Endothelialization & 6.90 & 19.51 & 0.00 & 0.148 \\
SMC\_inhibition & 2.66 & 14.47 & 0.00 & 0.057 \\
Anti\_inflammation & 12.97 & 27.04 & 0.00 & 0.278 \\
Thrombogenicity & 2.52 & 14.27 & 0.00 & 0.054 \\
Total\_Score & 3.96 & 10.52 & 0.00 & 0.085 \\
\bottomrule
\end{tabular}
\end{table}

\subsection{Inter-rater Reliability Analysis}

The intraclass correlation coefficient (ICC) reflects scoring consistency across different runs:

\begin{table}[H]
\centering
\caption{gemini-3-pro-preview - Intraclass Correlation Coefficient}
\begin{tabular}{lc}
\toprule
\textbf{Scoring Dimension} & \textbf{ICC} \\
\midrule
Mechanical\_Safety & 0.9298 \\
Swelling\_Performance & 0.8974 \\
Endothelialization & 0.9370 \\
SMC\_inhibition & 0.9425 \\
Anti\_inflammation & 0.7313 \\
Thrombogenicity & 0.9946 \\
Total\_Score & 0.9070 \\
\bottomrule
\end{tabular}
\end{table}

\subsection{Decision Consistency Analysis}


\begin{itemize}
    \item \textbf{Total Runs}: 11 runs
    \item \textbf{Most Common Winner}: Formula 4
    \item \textbf{Occurrence Count}: 11 times
    \item \textbf{Decision Consistency Rate}: 100.0\%
    \item \textbf{Information Entropy}: -0.000 (Max Entropy: 0.000)
\end{itemize}

Winner formula distribution:

\begin{table}[H]
\centering
\caption{gemini-3-pro-preview - Winner Formula Distribution}
\begin{tabular}{cc}
\toprule
\textbf{Formula} & \textbf{Count} \\
\midrule
Formula 4 & 11 \\
\bottomrule
\end{tabular}
\end{table}


\newpage


\section{Model Comparison Analysis}

\subsection{Overall Reliability Comparison}

The following table summarizes the reliability performance of the four LLM models across all metrics:

\begin{table}[H]
\centering
\caption{Model Reliability Metrics Comparison}
\begin{tabular}{lccc}
\toprule
\textbf{Model} & \textbf{Mean CV(\%)} & \textbf{Mean ICC} & \textbf{Decision Consistency(\%)} \\
\midrule
gpt-5 & 8.96 & 0.8126 & 90.9 \\
grok-4 & 9.72 & 0.5368 & 63.6 \\
claude-opus-4-5-20251101 & 6.42 & 0.8290 & 63.6 \\
gemini-3-pro-preview & 6.73 & 0.9056 & 100.0 \\
\bottomrule
\end{tabular}
\end{table}

\subsection{Reliability Ranking}

Based on comprehensive performance, we rank the models by reliability:

\begin{enumerate}
\item \textbf{gemini-3-pro-preview} (Overall Score: 93.1/100)
\begin{itemize}
    \item Scoring Consistency: 86.5
    \item Inter-rater Reliability: 90.6
    \item Decision Stability: 100.0
\end{itemize}

\item \textbf{gpt-5} (Overall Score: 85.4/100)
\begin{itemize}
    \item Scoring Consistency: 82.1
    \item Inter-rater Reliability: 81.3
    \item Decision Stability: 90.9
\end{itemize}

\item \textbf{claude-opus-4-5-20251101} (Overall Score: 76.5/100)
\begin{itemize}
    \item Scoring Consistency: 87.2
    \item Inter-rater Reliability: 82.9
    \item Decision Stability: 63.6
\end{itemize}

\item \textbf{grok-4} (Overall Score: 65.7/100)
\begin{itemize}
    \item Scoring Consistency: 80.6
    \item Inter-rater Reliability: 53.7
    \item Decision Stability: 63.6
\end{itemize}

\end{enumerate}

\newpage


\section{Conclusions and Discussion}

\subsection{Key Findings}

Based on systematic statistical analysis of four LLM models, we report the following key findings:

\subsubsection{gpt-5}

\begin{itemize}
    \item Very low scoring variability (CV=8.96\%), demonstrating excellent scoring consistency.
\end{itemize}
\begin{itemize}
    \item Very high inter-rater correlation (ICC=0.8126), indicating highly consistent scoring across runs.
\end{itemize}
\begin{itemize}
    \item Excellent decision stability with 90.9\% of runs selecting the same winner (Formula 5).
\end{itemize}
\subsubsection{grok-4}

\begin{itemize}
    \item Very low scoring variability (CV=9.72\%), demonstrating excellent scoring consistency.
\end{itemize}
\begin{itemize}
    \item Moderate inter-rater correlation (ICC=0.5368), indicating some variation in scoring patterns.
\end{itemize}
\begin{itemize}
    \item Good decision stability with 63.6\% of runs selecting the same winner (Formula 4).
\end{itemize}
\subsubsection{claude-opus-4-5-20251101}

\begin{itemize}
    \item Very low scoring variability (CV=6.42\%), demonstrating excellent scoring consistency.
\end{itemize}
\begin{itemize}
    \item Very high inter-rater correlation (ICC=0.8290), indicating highly consistent scoring across runs.
\end{itemize}
\begin{itemize}
    \item Good decision stability with 63.6\% of runs selecting the same winner (Formula 4).
\end{itemize}
\subsubsection{gemini-3-pro-preview}

\begin{itemize}
    \item Very low scoring variability (CV=6.73\%), demonstrating excellent scoring consistency.
\end{itemize}
\begin{itemize}
    \item Very high inter-rater correlation (ICC=0.9056), indicating highly consistent scoring across runs.
\end{itemize}
\begin{itemize}
    \item Excellent decision stability with 100.0\% of runs selecting the same winner (Formula 4).
\end{itemize}

\subsection{Statistical Reliability Assessment}

Based on the comprehensive analysis above, we draw the following statistical conclusions:

\begin{enumerate}
    \item \textbf{Scoring Reliability}: All models show good stability in most scoring dimensions, with CV普遍低于20\%, indicating that the same model produces consistent scores for the same material across different runs.

    \item \textbf{Inter-run Correlation}: The average ICC values for all models reach good or excellent levels, indicating that scoring patterns are highly similar across runs, and the model does not produce systematic biases due to random factors.

    \item \textbf{Decision Stability}: Most models show decision consistency rates exceeding 60\%, with some models reaching 80\% or above, demonstrating that the model can stably identify the optimal material candidate in repeated runs.
\end{enumerate}

\subsection{Research Conclusions}

\textbf{Core Conclusion}: This study fully confirms the hypothesis that \textit{multiple independent runs of the same AI model are statistically reliable}. This is demonstrated by:

\begin{itemize}
    \item Highly consistent scores for the same material across different runs (low coefficient of variation)
    \item Highly correlated scoring patterns across runs (high ICC values)
    \item Stable decision-making in final winner selection (high decision consistency rate)
\end{itemize}

This finding provides statistical evidence for applying LLM models in materials science research and decision-making, demonstrating the reliability and reproducibility of their output results.

\subsection{Limitations and Future Work}

\begin{itemize}
    \item This study analyzed only four LLM models; future work could extend to more models
    \item The number of runs (11) is relatively limited; increasing the number of runs may further improve statistical significance
    \item Future research could investigate the impact of different temperature parameters on reliability
    \item Could explore reliability performance in different task types (non-material selection)
\end{itemize}

\newpage


\appendix

\section{Detailed Data Tables}

\subsection{Scoring Statistics by Formula}

\subsubsection{gpt-5}

\begin{table}[H]
\centering
\caption{gpt-5 - Detailed Scoring Statistics}
\small
\begin{tabular}{llcccc}
\toprule
\textbf{Scoring Dimension} & \textbf{Formula} & \textbf{Mean} & \textbf{Std Dev} & \textbf{CV(\%)} & \textbf{Sample Size} \\
\midrule
Mechanical\_Safety & See detailed stats & See detailed stats & See detailed stats & See detailed stats & See detailed stats \\
Swelling\_Performance & See detailed stats & See detailed stats & See detailed stats & See detailed stats & See detailed stats \\
Endothelialization & See detailed stats & See detailed stats & See detailed stats & See detailed stats & See detailed stats \\
SMC\_inhibition & See detailed stats & See detailed stats & See detailed stats & See detailed stats & See detailed stats \\
Anti\_inflammation & See detailed stats & See detailed stats & See detailed stats & See detailed stats & See detailed stats \\
Thrombogenicity & See detailed stats & See detailed stats & See detailed stats & See detailed stats & See detailed stats \\
Total\_Score & See detailed stats & See detailed stats & See detailed stats & See detailed stats & See detailed stats \\
\bottomrule
\end{tabular}
\end{table}

Note: Detailed statistics for each formula are available in the JSON data file.

\subsubsection{grok-4}

\begin{table}[H]
\centering
\caption{grok-4 - Detailed Scoring Statistics}
\small
\begin{tabular}{llcccc}
\toprule
\textbf{Scoring Dimension} & \textbf{Formula} & \textbf{Mean} & \textbf{Std Dev} & \textbf{CV(\%)} & \textbf{Sample Size} \\
\midrule
Mechanical\_Safety & See detailed stats & See detailed stats & See detailed stats & See detailed stats & See detailed stats \\
Swelling\_Performance & See detailed stats & See detailed stats & See detailed stats & See detailed stats & See detailed stats \\
Endothelialization & See detailed stats & See detailed stats & See detailed stats & See detailed stats & See detailed stats \\
SMC\_inhibition & See detailed stats & See detailed stats & See detailed stats & See detailed stats & See detailed stats \\
Anti\_inflammation & See detailed stats & See detailed stats & See detailed stats & See detailed stats & See detailed stats \\
Thrombogenicity & See detailed stats & See detailed stats & See detailed stats & See detailed stats & See detailed stats \\
Total\_Score & See detailed stats & See detailed stats & See detailed stats & See detailed stats & See detailed stats \\
\bottomrule
\end{tabular}
\end{table}

Note: Detailed statistics for each formula are available in the JSON data file.

\subsubsection{claude-opus-4-5-20251101}

\begin{table}[H]
\centering
\caption{claude-opus-4-5-20251101 - Detailed Scoring Statistics}
\small
\begin{tabular}{llcccc}
\toprule
\textbf{Scoring Dimension} & \textbf{Formula} & \textbf{Mean} & \textbf{Std Dev} & \textbf{CV(\%)} & \textbf{Sample Size} \\
\midrule
Mechanical\_Safety & See detailed stats & See detailed stats & See detailed stats & See detailed stats & See detailed stats \\
Swelling\_Performance & See detailed stats & See detailed stats & See detailed stats & See detailed stats & See detailed stats \\
Endothelialization & See detailed stats & See detailed stats & See detailed stats & See detailed stats & See detailed stats \\
SMC\_inhibition & See detailed stats & See detailed stats & See detailed stats & See detailed stats & See detailed stats \\
Anti\_inflammation & See detailed stats & See detailed stats & See detailed stats & See detailed stats & See detailed stats \\
Thrombogenicity & See detailed stats & See detailed stats & See detailed stats & See detailed stats & See detailed stats \\
Total\_Score & See detailed stats & See detailed stats & See detailed stats & See detailed stats & See detailed stats \\
\bottomrule
\end{tabular}
\end{table}

Note: Detailed statistics for each formula are available in the JSON data file.

\subsubsection{gemini-3-pro-preview}

\begin{table}[H]
\centering
\caption{gemini-3-pro-preview - Detailed Scoring Statistics}
\small
\begin{tabular}{llcccc}
\toprule
\textbf{Scoring Dimension} & \textbf{Formula} & \textbf{Mean} & \textbf{Std Dev} & \textbf{CV(\%)} & \textbf{Sample Size} \\
\midrule
Mechanical\_Safety & See detailed stats & See detailed stats & See detailed stats & See detailed stats & See detailed stats \\
Swelling\_Performance & See detailed stats & See detailed stats & See detailed stats & See detailed stats & See detailed stats \\
Endothelialization & See detailed stats & See detailed stats & See detailed stats & See detailed stats & See detailed stats \\
SMC\_inhibition & See detailed stats & See detailed stats & See detailed stats & See detailed stats & See detailed stats \\
Anti\_inflammation & See detailed stats & See detailed stats & See detailed stats & See detailed stats & See detailed stats \\
Thrombogenicity & See detailed stats & See detailed stats & See detailed stats & See detailed stats & See detailed stats \\
Total\_Score & See detailed stats & See detailed stats & See detailed stats & See detailed stats & See detailed stats \\
\bottomrule
\end{tabular}
\end{table}

Note: Detailed statistics for each formula are available in the JSON data file.


\section*{References}

\begin{thebibliography}{9}

\bibitem{shrout1979}
Shrout PE, Fleiss JL.
\newblock Intraclass correlations: uses in assessing rater reliability.
\newblock \emph{Psychological Bulletin}. 1979;86(2):420-428.

\bibitem{mcgraw1996}
McGraw KO, Wong SP.
\newblock Forming inferences about some intraclass correlation coefficients.
\newblock \emph{Psychological Methods}. 1996;1(1):30-46.

\bibitem{landis1977}
Landis JR, Koch GG.
\newblock The measurement of observer agreement for categorical data.
\newblock \emph{Biometrics}. 1977;33(1):159-174.

\bibitem{koo2006}
Koo TK, Li JD.
\newblock A guideline of selecting and reporting intraclass correlation coefficients for reliability research.
\newblock \emph{Journal of Chiropractic Medicine}. 2016;19(3):342-349.

\end{thebibliography}

\newpage

\end{document}